\documentclass[12pt]{article}
\usepackage[utf8]{inputenc}
\usepackage[T1]{fontenc}
\usepackage{hyperref}
\title{Uduk Language Resources for K-12 Curriculum Development}
\author{Compiled by Claude Sonnet 5 (claude-sonnet-5) for Andr\'as Kornai}
\date{}
\begin{document}
\maketitle

\section{Introduction}
Uduk (ISO 639-3 udu; Glottocode uduk1239) is a Koman language of Sudan and South Sudan with about 22,000 speakers. It has three dialects: Chali and Bellila in the north, and Yabus in the south. This report rests on a small number of searches and is short. I found no dictionary, textbook or primer.

\section{Primary Educational Resources}
I found no Uduk textbooks or primers. Omniglot says Uduk was first written in the 1950s, that missionaries developed a standard way of writing it, and that it is used in some schools. \url{https://omniglot.com/writing/uduk.htm}

\section{Linguistic and Reference Materials}
\begin{itemize}
\item The Wikipedia article summarises the dialects, writing system and speaker numbers. \url{https://en.wikipedia.org/wiki/Uduk_language}
\item Don Killian's academic page at the University of Helsinki is the main source I found on Uduk grammar and phonology. \url{https://helsinki.academia.edu/DonKillian}
\end{itemize}

\section{Religious and Community Resources}
\begin{itemize}
\item Global Recordings lists Uduk audio Bible lessons and evangelistic programmes, with a link to a New Testament audio recording. \url{https://globalrecordings.net/en/language/udu}
\item Joshua Project has a profile of the Uduk people in Sudan. \url{https://www.joshuaproject.net/index.php/people_groups/15687/SU}
\end{itemize}

\section{Digital and Audio Resources}
The Global Recordings page above is the only audio source I found.

\section{Curriculum Development Context}
South Sudan has low literacy overall, 34.5\% according to a Bible Society fact sheet, which does not address Uduk specifically. \url{https://www.biblesociety.org.au/wp-content/uploads/2021/04/100525-Fact-Sheet-South-Sudan-Literacy.pdf}
Any curriculum would have to be built almost from scratch, working with the community and with the missionary orthography.

\end{document}